% ============================================================
% Abkürzungsverzeichnis
% ============================================================
% Verwendung im Text:
%   \ac{API}   → Beim 1. Vorkommen: Langform (Kurzform), danach: Kurzform
%   \acs{API}  → Immer nur Kurzform
%   \acl{API}  → Immer nur Langform
% ============================================================

\section*{Abkürzungsverzeichnis}
\addcontentsline{toc}{section}{Abkürzungsverzeichnis}
{\protect\thispagestyle{fancy}}
\addtocontents{toc}{\protect\vspace{-5pt}}

\vspace{0.25cm}

\begin{onehalfspace}
	\newcommand{\acroformat}[2]{\textbf{#1}\dotfill \nobreak\hspace{0.5em}#2}

	\makeatletter
	\renewcommand{\dotfill}{\leavevmode \cleaders \hb@xt@ .75em{\hss .\hss }\hfill \kern \z@}
	\makeatother

	\setstretch{0}

	\begin{acronym}
		% Syntax: \acro{KÜRZEL}[KÜRZEL]{\acroformat{}{Ausgeschriebene Bedeutung}}
		\acro{API}[API]{\acroformat{}{Application Programming Interface}}
		\acro{EU}[EU]{\acroformat{}{Europäische Union}}
		\acro{IT}[IT]{\acroformat{}{Informationstechnologie}}
		\acro{KI}[KI]{\acroformat{}{Künstliche Intelligenz}}
		\acro{KPI}[KPI]{\acroformat{}{Key Performance Indicator}}
		% Weitere Abkürzungen hier ergänzen …
	\end{acronym}
\end{onehalfspace}
